\section{成果二：证据条件化的网络路径优化}
\label{sec:network-method}
\subsection{由节点分数转向有前提的完整路径}
网络优化的目标为预算内的完整任务产物，而不是把高分节点排在前面。
把节点 $s$ 写为 $(H_s,R_s,A_s,c_s)$，分别表示实现资源指纹、必要前提集合、
完成后新增的能力事实及上下文字节成本。
初始事实集为原始CSV存在，目标事实为指标已完成。
路径 $P=(s_1,\ldots,s_l)$ 必须满足
\begin{align}
R_{s_i}&\subseteq F_{i-1},&F_i&=F_{i-1}\cup A_{s_i},\\
\sum_{s\in P}c_s&\leq B,&l&\leq5.
\end{align}
例如，任务同时要求按键去重、排除跨分割主体和删除空分数时，顺序为先去重、再排除、再删空，
指标节点依赖这些必要事实全部成立。
复合节点可一次提供三个清洗效果。
因此逐个取高分的前五个节点可能重复同一能力并遗漏必要阶段；
“节点各自看起来好”与“存在兼容的完整路径”不是相同条件。

\subsection{证据匹配、条件分组与搜索目标}
本文实现 evidence-conditioned path（ECP）。证据只来自开发集的独立产物探测。
匹配键为完整工作资源的 SHA-256 与可观察条件组，而不是节点名称。
解析节点按是否含CSV引号分组，指标节点按是否存在阈值边界分组，其余按全部条件合并。
条件标记由任务适配器提供，所有策略均可见；它们不是验收答案，也尚未由通用语义识别器自动提取。

对匹配组中的成功数 $u_{s,g}$ 与探测数 $n_{s,g}$，定义排序启发量
\begin{equation}
\widetilde q(s,g)=\frac{u_{s,g}+1}{n_{s,g}+2}.
\end{equation}
将已观测到组内失败的节点排除，再枚举预算内的兼容路径，选择
\begin{equation}
\begin{aligned}
J(P)&=\sum_{s\in P}-\log\widetilde q(s,g)\\
&\quad+0.05\frac{\sum_{s\in P}c_s}{B}+0.02|P|,\\
\widehat P&=\arg\min_{P\ \mathrm{feasible}}J(P).
\end{aligned}
\label{eq:ecp}
\end{equation}
系数、分组、预算档位及比较规则在测试输出产生前固定。
未探测节点取先验 $1/2$，仍可能被选中；这明确是探索性路由，不是把未知标成可信。
各节点相关、前提状态可能变化，故 $\widetilde q$ 的乘积不是校准的路径成功概率。
式\ref{eq:ecp}只是小网络可穷举的确定性目标，不主张新的最短路径理论或普遍最优。

\subsection{与判别器的连接及公平对照}
成果一的通过状态严格绑定已测输入；成果二将开发观察用作选择新任务的假设，
仍须执行新路径并通过新的独立验收。
两者不可混为“已知输入通过，所以未来输入都被认证”。
每个新产物同时核对完整保留行、行序、计数、混淆矩阵及准确率、精确率、召回率、F1。
任一必要检查不合格均记失败；没有路径的弃权也按未完成计入分母。
全部策略共用验收器，其阻断错误产物的作用不能算成ECP独有的学习能力。

十种策略共用候选、任务、预算和执行器：词法、MMR及DPP前五后接原生规划器；
项目原生元数据最短路径；预先配置的正确人工流程；全局证据过滤后最短路径；
全局证据风险路径；ECP；条件证据前五后规划；去版本绑定的ECP。
全局风险路径与ECP仅区别于是否按输入条件匹配，去版本策略则故意只按ID沿用证据。
原生路径基线直接调用项目已有 \texttt{skillnet.frontier.plan}，没有另写较弱替代品。
人工流程使用预先选定的标准正确组件，是已知工作流条件下的强参照；
其人工配置成本未量化，因此不能直接当作可扩展自动路由器。
MMR/DPP使用词法表示，未复现训练版SkillRouter或DSR，不能宣称超过原论文模型。
